%%%%%%%%%%%%%%%%%%%%%%%%%%%%%%%%%%%%%%%%%%%%%%%%%%%%%%%%%%%%%%%%%%%%%%
% Overleaf Example: A quick guide to LaTeX
%
% Original Source: Dave Richeson (divisbyzero.com), Dickinson College
% Modified By: Paul Gessler, Overleaf (overleaf.com)
% 
% A one-size-fits-all LaTeX cheat sheet. Kept to two pages, so it 
% can be printed (double-sided) on one piece of paper
% 
% Feel free to distribute this example, but please keep the referral
% to divisbyzero.com
% 
% Guidance on the use of the Overleaf logos can be found here:
% https://www.overleaf.com/for/partners/logos 
%%%%%%%%%%%%%%%%%%%%%%%%%%%%%%%%%%%%%%%%%%%%%%%%%%%%%%%%%%%%%%%%%%%%%%
% How to use Overleaf: 
%
% You edit the source code here on the left, and the preview on the
% right shows you the result within a few seconds.
%
% Bookmark this page and share the URL with your co-authors. They can
% edit at the same time!
%
% You can upload figures, bibliographies, custom classes and
% styles using the files menu.
%
% If you're new to LaTeX, the wikibook is a great place to start:
% http://en.wikibooks.org/wiki/LaTeX
%
%%%%%%%%%%%%%%%%%%%%%%%%%%%%%%%%%%%%%%%%%%%%%%%%%%%%%%%%%%%%%%%%%%%%%%

\documentclass[10pt,landscape,letterpaper]{article}
\usepackage{amssymb,amsmath,amsthm,amsfonts}
\usepackage[table,xcdraw]{xcolor}

\usepackage{multicol,multirow}
\usepackage{spverbatim}
\usepackage{graphicx}
\usepackage{ifthen}
\usepackage[landscape]{geometry}
\usepackage[colorlinks=true,urlcolor=olgreen]{hyperref}
\usepackage{booktabs}
\usepackage{fontspec}
\setmainfont[Ligatures=TeX]{TeX Gyre Pagella}
\setsansfont{Fira Sans}
\setmonofont{Inconsolata}
\usepackage{unicode-math}
\setmathfont{TeX Gyre Pagella Math}
\usepackage{microtype}
\usepackage{xcolor}
% dirty fix for microtype
\makeatletter
\def\MT@is@composite#1#2\relax{%
  \ifx\\#2\\\else
    \expandafter\def\expandafter\MT@char\expandafter{\csname\expandafter
                    \string\csname\MT@encoding\endcsname
                    \MT@detokenize@n{#1}-\MT@detokenize@n{#2}\endcsname}%
    % 3 lines added:
    \ifx\UnicodeEncodingName\@undefined\else
      \expandafter\expandafter\expandafter\MT@is@uni@comp\MT@char\iffontchar\else\fi\relax
    \fi
    \expandafter\expandafter\expandafter\MT@is@letter\MT@char\relax\relax
    \ifnum\MT@char@ < \z@
      \ifMT@xunicode
        \edef\MT@char{\MT@exp@two@c\MT@strip@prefix\meaning\MT@char>\relax}%
          \expandafter\MT@exp@two@c\expandafter\MT@is@charx\expandafter
            \MT@char\MT@charxstring\relax\relax\relax\relax\relax
      \fi
    \fi
  \fi
}
% new:
\def\MT@is@uni@comp#1\iffontchar#2\else#3\fi\relax{%
  \ifx\\#2\\\else\edef\MT@char{\iffontchar#2\fi}\fi
}
\makeatother

\ifthenelse{\lengthtest { \paperwidth = 11in}}
    { \geometry{margin=0.4in} }
	{\ifthenelse{ \lengthtest{ \paperwidth = 297mm}}
		{\geometry{top=1cm,left=1cm,right=1cm,bottom=1cm} }
		{\geometry{top=1cm,left=1cm,right=1cm,bottom=1cm} }
	}
\pagestyle{empty}
\makeatletter
\renewcommand{\section}{\@startsection{section}{1}{0mm}%
                                {-1ex plus -.5ex minus -.2ex}%
                                {0.5ex plus .2ex}%x
                                {\sffamily\large}}
\renewcommand{\subsection}{\@startsection{subsection}{2}{0mm}%
                                {-1explus -.5ex minus -.2ex}%
                                {0.5ex plus .2ex}%
                                {\sffamily\normalsize\itshape}}
\renewcommand{\subsubsection}{\@startsection{subsubsection}{3}{0mm}%
                                {-1ex plus -.5ex minus -.2ex}%
                                {1ex plus .2ex}%
                                {\normalfont\small\itshape}}
\makeatother
\setcounter{secnumdepth}{0}
\setlength{\parindent}{0pt}
\setlength{\parskip}{0pt plus 0.5ex}
% -----------------------------------------------------------------------

\usepackage{academicons}

\usepackage{draftwatermark,afterpage,xcolor}
\definecolor{olgreen}{HTML}{4f9c45}
\SetWatermarkText{}
\SetWatermarkFontSize{0.1\paperheight}
\SetWatermarkAngle{0}
\SetWatermarkColor[HTML]{EDF5EC}
%\SetWatermarkLightness{0.95}
\SetWatermarkHorCenter{0.35\paperwidth}

\begin{document}
\footnotesize
%\raggedright

\begin{center}
  {\huge\sffamily\bfseries Aide Mémoire de la grammaire Turque (A1-A2)} 
\end{center} 
\setlength{\premulticols}{0pt}
\setlength{\postmulticols}{0pt}
\setlength{\multicolsep}{1pt}
\setlength{\columnsep}{1.8em}
\begin{multicols}{3}




\section{\underline{Structure de la langue}}

 La langue turque est une langue de type SOV (Sujet – Objet – Verbe). 
Elle se caractérise par un fonctionnement agglutinant, dans lequel les 
mots sont formés et enrichis par l’ajout de suffixes successifs. Le turc
 ne distingue pas le masculin du 
féminin.  

\section{\underline{Harmonies vocaliques}}

La langue turque obéit à des règles d’harmonie vocalique : la 
forme des suffixes varie en fonction des voyelles présentes dans la 
dernière syllabe du radical.

\begin{tabular}{|l|ll|ll|}
\hline
\multicolumn{1}{|c|}{}                          & \multicolumn{2}{l|}{Voyelle suivante}                                                                          & \multicolumn{2}{l|}{Exemples}                                                                                    \\ \cline{2-5} 
\multicolumn{1}{|c|}{\multirow{-2}{*}{Voyelle}} & \multicolumn{1}{l|}{Type i}                                      & Type e                                      & \multicolumn{1}{l|}{Type i}                                                   & Type e                           \\ \hline \hline
a                                               & \multicolumn{1}{l|}{\cellcolor[HTML]{C0C0C0}}                    & \cellcolor[HTML]{EFEFEF}                    & \multicolumn{1}{l|}{\cellcolor[HTML]{C0C0C0}{\color[HTML]{333333} kitap mı?}} & \cellcolor[HTML]{EFEFEF}kitaplar \\ \cline{1-1} \cline{4-5} 
ı                                               & \multicolumn{1}{l|}{\multirow{-2}{*}{\cellcolor[HTML]{C0C0C0}ı}} & \cellcolor[HTML]{EFEFEF}                    & \multicolumn{1}{l|}{\cellcolor[HTML]{C0C0C0}{\color[HTML]{333333} kapı mı?}}  & \cellcolor[HTML]{EFEFEF}kapılar  \\ \cline{1-2} \cline{4-5} 
o                                               & \multicolumn{1}{l|}{}                                            & \cellcolor[HTML]{EFEFEF}                    & \multicolumn{1}{l|}{top mu?}                                                  & \cellcolor[HTML]{EFEFEF}toplar   \\ \cline{1-1} \cline{4-5} 
u                                               & \multicolumn{1}{l|}{\multirow{-2}{*}{u}}                         & \multirow{-4}{*}{\cellcolor[HTML]{EFEFEF}a} & \multicolumn{1}{l|}{okul mu?}                                                 & \cellcolor[HTML]{EFEFEF}okullar  \\ \hline \hline
e                                               & \multicolumn{1}{l|}{\cellcolor[HTML]{C0C0C0}}                    & \cellcolor[HTML]{EFEFEF}                    & \multicolumn{1}{l|}{\cellcolor[HTML]{C0C0C0}gece mi?}                         & \cellcolor[HTML]{EFEFEF}geceler  \\ \cline{1-1} \cline{4-5} 
i                                               & \multicolumn{1}{l|}{\multirow{-2}{*}{\cellcolor[HTML]{C0C0C0}i}} & \cellcolor[HTML]{EFEFEF}                    & \multicolumn{1}{l|}{\cellcolor[HTML]{C0C0C0}tatil mi?}                        & \cellcolor[HTML]{EFEFEF}tatiller \\ \cline{1-2} \cline{4-5} 
ö                                               & \multicolumn{1}{l|}{}                                            & \cellcolor[HTML]{EFEFEF}                    & \multicolumn{1}{l|}{göl mü?}                                                  & \cellcolor[HTML]{EFEFEF}göller   \\ \cline{1-1} \cline{4-5} 
ü                                               & \multicolumn{1}{l|}{\multirow{-2}{*}{u}}                         & \multirow{-4}{*}{\cellcolor[HTML]{EFEFEF}e} & \multicolumn{1}{l|}{gün mü?}                                                  & \cellcolor[HTML]{EFEFEF}günler   \\ \hline
\end{tabular}

\section{\underline{Régles sur les consonnes}}

fıstıkçı şahap
p t 


\section{\underline{Les cas grammaticaux}}

Un cas est un suffixe ajouté à un nom (ou pronom) pour indiquer sa fonction dans la phrase, comme le rôle du sujet, du complément d’objet ou du lieu.

{
\setlength{\tabcolsep}{1pt}
\begin{tabular}{|
>{\columncolor[HTML]{EFEFEF}}l |p{3cm}|
>{\columncolor[HTML]{EFEFEF}}l |p{3cm}|}
\hline
Cas       & Définition                              & Suffixe & Exemple               \\ \hline \hline
Nom. & sujet ou complément \shortstack{indéfini d'un verbe} &         & kedi bir fare yiyor  \shortstack{\textit{le chat mange une souris}}              \\ \hline
Acc. & complément d'objet \shortstack{direct défini d'un verbe}          & -(y)i  & kedi fareyi yiyor  \shortstack{\textit{le chat mange la souris}}  \\ \hline
Gén.   & possession      & -(n)in & kedinin mutfağı \shortstack{\textit{le chat de la cuisine}}    \\ \hline
Dir.  & direction                  & -(y)e   & kedi mutfağa gidiyor \shortstack{\textit{le chat va dans la cuisine}} \\ \hline
Loc.   & lieu où une personne ou un objet se situe & -(d)e  & kedi mutfakta yiyor \shortstack{\textit{le chat mange dans cuisine}} \\ \hline
Abl.   & provenance                    & -(d)en  & kedi mutfaktan geliyor \textit{le chat vient de la cuisine}\\ \hline
\end{tabular}
}

\section{\underline{La possession en détail}}

Couple Génitif + Accusatif et sinon pronom +  

\section{\underline{Les principaux suffixes}}

siz, li, le, lik, ci


\section{\underline{Les adjectifs}}

Avant le nom pour épithéte pour qualifier le nom et après pour définir le nom et règle avec bir



\columnbreak

\section{\underline{Les pronoms personnels}}
{
\setlength{\tabcolsep}{3pt}
\begin{tabular}{|c|c|c|c|c|c|c|}
\hline
\cellcolor[HTML]{C0C0C0} & sujet & acc.   & gén.    & dir.   & loc     & abl.    \\ \hline \hline
je                       & ben   & beni   & benin   & bana   & bende   & benden  \\ \hline
tu                       & sen   & seni   & senin   & sana   & sende   & senden  \\ \hline
il/elle                  & o     & onu    & onun    & ona    & onda    & ondan   \\ \hline \hline
nous                     & biz   & bizi   & bizim   & bize   & bizde   & bizden  \\ \hline
vous                     & siz   & sizi   & sizin   & size   & sizde   & sizden  \\ \hline
ils/elles                & onlar & onları & onların & onlara & onlarda & olardan \\ \hline
\end{tabular}
}

\section{\underline{La conjugaison du verbe "être"}}

Au présent,...

Au passé,...


\section{\underline{La conjugaison du verbe "avoir"}}

Il n'existe pas directement de verbe "avoir" dans la lanque turque,  mais seulement deux mots qui var et yok qui traduise l'e

\section{\underline{La conjugaison des autres verbes}}

Les verbes sont donnés à l'infinitif avec le suffixe -mek. Le radical s'obtient en supprimant ce suffixe : gelmek \rightarrow{} gel,  yapmak \rightarrow{} yap. 

\subsection{Le présent continu}

Le présent continu sert principalement à exprimer une action en train de se dérouler au moment où l’on parle. C’est l’équivalent direct du français “être en train de…”.

\begin{tabular}{|lll|}
\hline
\multicolumn{3}{|c|}{Radical + \textcolor{red}{iyor} + \textcolor{blue}{suffixe personne}}                                                                    \\ \hline
\multicolumn{1}{|l|}{Affirmatif}   & \multicolumn{1}{l|}{Négatif}       & Interrogatif     \\ \hline
\multicolumn{1}{|l|}{gel\textcolor{red}{iyor}\textcolor{blue}{um}}    & \multicolumn{1}{l|}{gelm\textcolor{red}{iyor}\textcolor{blue}{um}}    & gel\textcolor{red}{iyor} muy\textcolor{blue}{um}?   \\ \hline
\multicolumn{1}{|l|}{gel\textcolor{red}{iyor}\textcolor{blue}{sun}}   & \multicolumn{1}{l|}{gelm\textcolor{red}{iyor}\textcolor{blue}{sun}}   & gel\textcolor{red}{iyor} mu\textcolor{blue}{sun}?   \\ \hline
\multicolumn{1}{|l|}{gel\textcolor{red}{iyor}} & \multicolumn{1}{l|}{gelm\textcolor{red}{iyor}}      & gel\textcolor{red}{iyor} mu?      \\ \hline\hline
\multicolumn{1}{|l|}{gel\textcolor{red}{iyor}\textcolor{blue}{uz}}    & \multicolumn{1}{l|}{gelm\textcolor{red}{iyor}\textcolor{blue}{uz}}    & gel\textcolor{red}{iyor} muy\textcolor{blue}{uz}?   \\ \hline
\multicolumn{1}{|l|}{gel\textcolor{red}{iyor}\textcolor{blue}{sunuz}} & \multicolumn{1}{l|}{gelm\textcolor{red}{iyor}\textcolor{blue}{sunuz}} & gel\textcolor{red}{iyor} mu\textcolor{blue}{sunuz}? \\ \hline
\multicolumn{1}{|l|}{gel\textcolor{red}{iyor}\textcolor{blue}{lar}}   & \multicolumn{1}{l|}{gelm\textcolor{red}{iyor}\textcolor{blue}{lar}}   & gel\textcolor{red}{iyor}\textcolor{blue}{lar} mu?   \\ \hline
\end{tabular}

Yemek et Etmek sont irréguliers.





\subsection{Le passé constaté}

\begin{tabular}{|lll|}
\hline
\multicolumn{3}{|c|}{Radical + \textcolor{red}{di} + \textcolor{blue}{suffixe personne}}                                                                    \\ \hline
\multicolumn{1}{|l|}{Affirmatif}   & \multicolumn{1}{l|}{Négatif}       & Interrogatif     \\ \hline
\multicolumn{1}{|l|}{gel\textcolor{red}{di}\textcolor{blue}{m}}    & \multicolumn{1}{l|}{gelme\textcolor{red}{di}\textcolor{blue}{m}}    & gel\textcolor{red}{di}\textcolor{blue}{m} mi?   \\ \hline
\multicolumn{1}{|l|}{gel\textcolor{red}{di}\textcolor{blue}{n}}   & \multicolumn{1}{l|}{gelme\textcolor{red}{di}\textcolor{blue}{n}}   & gel\textcolor{red}{di}\textcolor{blue}{n} mi?   \\ \hline
\multicolumn{1}{|l|}{gel\textcolor{red}{di}} & \multicolumn{1}{l|}{gelme\textcolor{red}{di}}      & gel\textcolor{red}{di} mi?      \\ \hline
\multicolumn{1}{|l|}{gel\textcolor{red}{di}\textcolor{blue}{k}}    & \multicolumn{1}{l|}{gelme\textcolor{red}{di}\textcolor{blue}{k}}    & gel\textcolor{red}{di}\textcolor{blue}{k} mi?   \\ \hline
\multicolumn{1}{|l|}{gel\textcolor{red}{di}\textcolor{blue}{niz}} & \multicolumn{1}{l|}{gelme\textcolor{red}{di}\textcolor{blue}{niz}} & gel\textcolor{red}{di}\textcolor{blue}{niz} mi? \\ \hline
\multicolumn{1}{|l|}{gel\textcolor{red}{di}\textcolor{blue}{ler}}   & \multicolumn{1}{l|}{gelme\textcolor{red}{di}\textcolor{blue}{ler}}   & gel\textcolor{red}{di}\textcolor{blue}{ler} mi?   \\ \hline
\end{tabular}


\subsection{Le futur}

\begin{tabular}{|lll|}
\hline
\multicolumn{3}{|c|}{Radical + \textcolor{red}{ecek} + \textcolor{blue}{suffixe personne}}                                                                    \\ \hline
\multicolumn{1}{|l|}{Affirmatif}   & \multicolumn{1}{l|}{Négatif}       & Interrogatif     \\ \hline \hline
\multicolumn{1}{|l|}{gel\textcolor{red}{ecek}\textcolor{blue}{im}}    & \multicolumn{1}{l|}{gelmey\textcolor{red}{eceğ}\textcolor{blue}{um}}    & gel\textcolor{red}{eceğ} muy\textcolor{blue}{um}?   \\ \hline
\multicolumn{1}{|l|}{gel\textcolor{red}{ecek}\textcolor{blue}{sin}}   & \multicolumn{1}{l|}{gelmey\textcolor{red}{ecek}\textcolor{blue}{sun}}   & gel\textcolor{red}{ecek} mu\textcolor{blue}{sun}?   \\ \hline
\multicolumn{1}{|l|}{gel\textcolor{red}{ecek}} & \multicolumn{1}{l|}{gelmey\textcolor{red}{ecek}}      & gel\textcolor{red}{ecek} mu?      \\ \hline
\multicolumn{1}{|l|}{gel\textcolor{red}{ecek}\textcolor{blue}{iz}}    & \multicolumn{1}{l|}{gelmey\textcolor{red}{ecek}\textcolor{blue}{uz}}    & gel\textcolor{red}{ecek} muy\textcolor{blue}{uz}?   \\ \hline
\multicolumn{1}{|l|}{gel\textcolor{red}{ecek}\textcolor{blue}{siniz}} & \multicolumn{1}{l|}{gelmey\textcolor{red}{ecek}\textcolor{blue}{sunuz}} & gel\textcolor{red}{ecek} mu\textcolor{blue}{sunuz}? \\ \hline
\multicolumn{1}{|l|}{gel\textcolor{red}{ecek}\textcolor{blue}{ler}}   & \multicolumn{1}{l|}{gelmey\textcolor{red}{ecek}\textcolor{blue}{lar}}   & gel\textcolor{red}{ecek}\textcolor{blue}{lar} mu?   \\ \hline
\end{tabular}


\subsection{L'impératif et l'optatif}

Il est principalement utilisé avec Sen et Siz :

\begin{tabular}{|ll|}
\hline
\multicolumn{2}{|l|}{Radical + \textcolor{blue}{suffixe personne}} \\ \hline
\multicolumn{1}{|l|}{Sen}         & gel          \\ \hline
\multicolumn{1}{|l|}{Siz}         & gelin        \\ \hline
\end{tabular}

\subsection{La capacité de faire une action}

radical + ebil et abil + present simple + suffixe personne

\bigskip

\fbox{Auteur : François Marouchi Sémécurbe   }

\end{multicols}

%\SetWatermarkHorCenter{0.65\paperwidth}
\end{document}
